\documentclass[UTF8,12pt]{ctexart}
\usepackage{amsmath,amssymb,geometry,bm,graphicx,fontspec,amssymb,amsthm}
\usepackage[mathscr]{euscript}


\usepackage{multirow}
\usepackage{tabularx}
\usepackage{tabularray}

\usepackage[colorlinks,
linkcolor=black,
anchorcolor=blue,
citecolor=green
]{hyperref} % 目录中的超链接
\usepackage{cleveref}

%数学定理
\newtheorem{Def}{定义}[section]
\newtheorem{Theo}{定理}[section]
\newtheorem{Lemm}{引理}[section]
\newtheorem{Prop}{命题}[section]
\newtheorem{Assu}{假设}[section]
\newtheorem{Axiom}{Axiom}

\newenvironment{Dr}{%
	\begin{list}{}%
		{
			\setlength{\leftmargin}{2em}
			\setlength{\labelwidth}{2em}
			\setlength{\labelsep}{0pt}
			\setlength{\itemindent}{0pt}
			\setlength{\listparindent}{0pt}
			\setlength{\parsep}{0pt}
			\setlength{\topsep}{0pt}
		}
		\item[\textbf{借：}]
	}{%
	\end{list}
}
\newenvironment{Cr}{%
	\begin{list}{}%
		{
			\setlength{\leftmargin}{2em}
			\setlength{\labelwidth}{2em}
			\setlength{\labelsep}{0pt}
			\setlength{\itemindent}{0pt}
			\setlength{\listparindent}{0pt}
			\setlength{\parsep}{0pt}
			\setlength{\topsep}{0pt}
		}
		\item[\textbf{贷：}]
	}{%
	\end{list}
}

\numberwithin{equation}{section} % 按章节进行排序与编号
\numberwithin{figure}{section}
\numberwithin{table}{section}

\usepackage{draftwatermark} % 所有页加水印
\SetWatermarkText{EconNerd} % 设置水印内容
\SetWatermarkLightness{0.99} % 设置水印透明度 0-1
\SetWatermarkScale{1} % 设置水印大小    

\title{题目} % 文档相关信息
\author{EconNerd}
\date{\today}
\geometry{scale=0.8}

\begin{document}
	\maketitle
	\tableofcontents
	\newpage
	
	\section{会计}
	\subsection{第一章 总论}
	\subsubsection{单选题}
	1. b
	
	2. d
	
	3. a（b）
	
	a选项，财务报告目标的受托责任观主要源于公司财产所有权和经营权的分离，要求财务报告应当有效反映受托者受托管理委托人财产责任的履行情况；财务报告目标的决策有用观主要源于资本市场的发展，要求财务报告需要向投资者提供与其投资决策相关的信息。c现金流量表为收付实现制。d政府会计主体的预算会计为收付实现制，财务会计为权责发生制
	
	4. b
	
	5. c
	
	6. d
	
	7. a
	
	8. a
	
	9. c
	
	10. a
	
	\subsubsection{多选题}
	1. abde（ade）
	
	诚信是注册会计师行业存在和发展的基石，在职业道德基本原则中居于首要地位
	
	2. acd（ac）
	
	些影响企业财务状况和经营成果的因素，如企业经营战略、研发能力、
	市场竞争力等，难以用货币来计量，企业可以在财务报告中补充披露有关非财务
	信息来弥补上述缺陷。
	
	3. bc
	
	4. abeg（abdeg）
	
	5. ab
	
	6. abc
	
	7. bcdef（bcef）
	
	于可预见的潜在的消极影响，应当以消极影响的严重程度和发生的可
	能性为评估标准。
	
	8. abcde（abde）
	
	质重于形式为会计信息质量要求。
	
	9. abcde（bce）
	
	\subsection{第二章 存货}
	\subsubsection{单选题}
	1. b
	
	2. c
	
	3. d（b）
	
	4. b（a）
	
	5. a
	
	6. d
	
	7. c
	
	8. c
	
	9. d
	
	10. b（a）存货跌价准备的转回
	
	\subsubsection{多选题}
	1. bcd（bce）
	
	委托外单位加工应税消费品(非金银首饰)应当区分收回委托加工物资
	后是否直接对外出售：(1)若收回后以不高于受托方计税价格直接对外销售，应
	当将消费税计入存货成本。(2)若收回后用于连续生产应税消费品，则消费税应
	当记入“应交税费——应交消费税”科目的借方，不计入存货成本。
	
	2. cd（ac）
	
	不允许后进后出，这个属于会计差错
	
	3. bc（cd）
	
	确定可变现净值的前提是持续经营假设
	
	4. abcd（abc）
	
	5. abc（acd）
	
	考虑管理不善的时候需要转出增值税
	
	6. abcdef（ade）
	
	企业销售自用生产设备不属于日常活动，若未满足划分为持有待售类别
	的条件，应当作为固定资产列报，若满足划分为持有待售类别的条件，应当作为
	持有待售资产列报。
	
	选项C,企业购入用于建造固定资产(办公楼)的物资，应当记入“工程物资”科目，
	列报为在建工程项目。
	
	选项F,企业在日常活动中持有，最终目的用于出售的数据资源应当确认为存货，
	在生产经营活动中使用的数据资源应当确认为无形资产。
	
	\subsection{第三章 固定资产}
	
	\subsubsection{单选题}
	1. a
	
	2. c
	
	3. b
	
	4. c
	
	5. d 看看c（d）
	
	有待售资产属于流动资产，不计提折旧。定资产的使用寿命、预计净残值和折旧方法的改变应作为会计估计变
	更进行处理。
	
	6. c
	
	7. d（b）
	
	8. d
	
	\subsubsection{多选题}
	1. cd（bcd）需要计算符合资本化的利息
	
	2. ac（abd）
	
	业提取的安全生产费应当记入“专项储备”科目，企业使用提取的安全生产费形成固定资产的，按形成固定资产的成本冲减专项储备，并确认相同金
	额的累计折旧。
	选项D,一般工商企业的固定资产发生的报废清理费用不属于弃置费用，应当在
	发生时作为固定资产处置费用处理。
	
	3. bc(bcd)
	
	×25年1月1日因技术进步甲企业应调减预计负债和固定资产的金额=
	522.3×(1+6?2600×(P/F,6?9)=73.68(万元)。2×25年1月1日固
	定资产的账面价值=50522.3-50522.3/30-73.68=48764.54(万元)。
	
	4. ad(ac)
	
	公司取得固定资产时，应按照所形成固定资产的成本冲减
	专项储备，并确认相同金额的累计折旧，因此2×24年12月31日甲公司“专项储
	备”科目的账面余额=200+150-300=50(万元),“累计折旧”科目的账面余额为
	300万元。
	
	5. abd(ad)
	
	选项B,生产设备与存货的生产加工相关，其日常维修费应当计入存货成
	本(制造费用)。
	
	6. abcd
	
	7. ac
	
	8. cd(ac)选项A,固定资产盘亏净损失，应当计入营业外支出。
	
	
	\subsection{第四章 无形资产}
	
	\subsubsection{单选题}
	1. d
	
	2. b 看下其他的解析
	
	3. c
	
	4. d（看看其他解析）
	
	5. d（a）考虑新的预计净残值
	
	6. b（d）
	
	7. c
	
	8. a
	
	\subsubsection{多选题}
	1. ac
	
	2. bcd
	
	3. abde
	
	4. cd（ad）
	
	5. ad（acd）
	
	6. cdf
	
	7. ac（abcd）
	
	8. bcd（cd）
	
	
	\subsection{第五章 投资性房地产}
	
	\subsubsection{单选题}
	1. a（c）
	
	】选项A,投资性房地产的范围包括已出租的土地使用权、持有并准备增值
	后转让的土地使用权和已出租的建筑物。因此非房地产开发企业持有并准备增值
	后转让的办公楼不属于投资性房地产，应当将其作为固定资产确认。
	
	2. d
	
	3. a
	
	4. a
	
	5. c（b）
	
	6. c
	
	7. c（b）产负债表日公允价值大于账面价值的差额计入公允价值变动损益。
	
	8. b（c） 公允价值变动损益不影响营业利润
	
	\subsubsection{多选题}
	1. acd（ad）
	
	地产开发企业用于开发对外出售写字楼的土地使用权应作为存货核算。
	
	2. abcde
	
	3. ac（ad）
	
	甲公司应当将该办公楼(投资性房地产)的费用化支出计入其他业务成本。
	
	4. ab（abd）
	
	5. bde（bce）
	
	选项D,自用建筑物停止自用改为出租，其转换日为租赁期开始日。只有空置建
	筑物或在建建筑物，如董事会或类似机构作出书面决议，明确表明将其用于经营
	出租且持有意图短期内不再发生变化，才能将其于作出书面决议时确认为投资性
	房地产。
	
	6. bc
	
	7. bc
	
	8. abcd（ab）
	
\end{document}